\documentclass[a4paper,11pt,oneside]{memoir}

% Castellano
\usepackage[spanish,es-tabla]{babel}
\selectlanguage{spanish}
\usepackage[utf8]{inputenc}
\usepackage{placeins}

\RequirePackage{booktabs}
\RequirePackage[table]{xcolor}
\RequirePackage{xtab}
\RequirePackage{multirow}

% Links
\usepackage[colorlinks]{hyperref}
\hypersetup{
	allcolors = {red}
}

% Ecuaciones
\usepackage{amsmath}

% Rutas de fichero / paquete
\newcommand{\ruta}[1]{{\sffamily #1}}

% Párrafos
\nonzeroparskip


% Imagenes
\usepackage{graphicx}
\newcommand{\imagen}[2]{
	\begin{figure}[!h]
		\centering
		\includegraphics[width=0.5\textwidth]{#1}
		\caption{#2}\label{fig:#1}
	\end{figure}
	\FloatBarrier
}

\newcommand{\imagenflotante}[2]{
	\begin{figure}%[!h]
		\centering
		\includegraphics[width=0.9\textwidth]{#1}
		\caption{#2}\label{fig:#1}
	\end{figure}
}



% El comando \figura nos permite insertar figuras comodamente, y utilizando
% siempre el mismo formato. Los parametros son:
% 1 -> Porcentaje del ancho de página que ocupará la figura (de 0 a 1)
% 2 --> Fichero de la imagen
% 3 --> Texto a pie de imagen
% 4 --> Etiqueta (label) para referencias
% 5 --> Opciones que queramos pasarle al \includegraphics
% 6 --> Opciones de posicionamiento a pasarle a \begin{figure}
\newcommand{\figuraConPosicion}[6]{%
  \setlength{\anchoFloat}{#1\textwidth}%
  \addtolength{\anchoFloat}{-4\fboxsep}%
  \setlength{\anchoFigura}{\anchoFloat}%
  \begin{figure}[#6]
    \begin{center}%
      \Ovalbox{%
        \begin{minipage}{\anchoFloat}%
          \begin{center}%
            \includegraphics[width=\anchoFigura,#5]{#2}%
            \caption{#3}%
            \label{#4}%
          \end{center}%
        \end{minipage}
      }%
    \end{center}%
  \end{figure}%
}

%
% Comando para incluir imágenes en formato apaisado (sin marco).
\newcommand{\figuraApaisadaSinMarco}[5]{%
  \begin{figure}%
    \begin{center}%
    \includegraphics[angle=90,height=#1\textheight,#5]{#2}%
    \caption{#3}%
    \label{#4}%
    \end{center}%
  \end{figure}%
}
% Para las tablas
\newcommand{\otoprule}{\midrule [\heavyrulewidth]}
%
% Nuevo comando para tablas pequeñas (menos de una página).
\newcommand{\tablaSmall}[5]{%
 \begin{table}
  \begin{center}
   \rowcolors {2}{gray!35}{}
   \begin{tabular}{#2}
    \toprule
    #4
    \otoprule
    #5
    \bottomrule
   \end{tabular}
   \caption{#1}
   \label{tabla:#3}
  \end{center}
 \end{table}
}

%
% Nuevo comando para tablas pequeñas (menos de una página).
\newcommand{\tablaSmallSinColores}[5]{%
 \begin{table}[H]
  \begin{center}
   \begin{tabular}{#2}
    \toprule
    #4
    \otoprule
    #5
    \bottomrule
   \end{tabular}
   \caption{#1}
   \label{tabla:#3}
  \end{center}
 \end{table}
}

\newcommand{\tablaApaisadaSmall}[5]{%
\begin{landscape}
  \begin{table}
   \begin{center}
    \rowcolors {2}{gray!35}{}
    \begin{tabular}{#2}
     \toprule
     #4
     \otoprule
     #5
     \bottomrule
    \end{tabular}
    \caption{#1}
    \label{tabla:#3}
   \end{center}
  \end{table}
\end{landscape}
}

%
% Nuevo comando para tablas grandes con cabecera y filas alternas coloreadas en gris.
\newcommand{\tabla}[6]{%
  \begin{center}
    \tablefirsthead{
      \toprule
      #5
      \otoprule
    }
    \tablehead{
      \multicolumn{#3}{l}{\small\sl continúa desde la página anterior}\\
      \toprule
      #5
      \otoprule
    }
    \tabletail{
      \hline
      \multicolumn{#3}{r}{\small\sl continúa en la página siguiente}\\
    }
    \tablelasttail{
      \hline
    }
    \bottomcaption{#1}
    \rowcolors {2}{gray!35}{}
    \begin{xtabular}{#2}
      #6
      \bottomrule
    \end{xtabular}
    \label{tabla:#4}
  \end{center}
}

%
% Nuevo comando para tablas grandes con cabecera.
\newcommand{\tablaSinColores}[6]{%
  \begin{center}
    \tablefirsthead{
      \toprule
      #5
      \otoprule
    }
    \tablehead{
      \multicolumn{#3}{l}{\small\sl continúa desde la página anterior}\\
      \toprule
      #5
      \otoprule
    }
    \tabletail{
      \hline
      \multicolumn{#3}{r}{\small\sl continúa en la página siguiente}\\
    }
    \tablelasttail{
      \hline
    }
    \bottomcaption{#1}
    \begin{xtabular}{#2}
      #6
      \bottomrule
    \end{xtabular}
    \label{tabla:#4}
  \end{center}
}

%
% Nuevo comando para tablas grandes sin cabecera.
\newcommand{\tablaSinCabecera}[5]{%
  \begin{center}
    \tablefirsthead{
      \toprule
    }
    \tablehead{
      \multicolumn{#3}{l}{\small\sl continúa desde la página anterior}\\
      \hline
    }
    \tabletail{
      \hline
      \multicolumn{#3}{r}{\small\sl continúa en la página siguiente}\\
    }
    \tablelasttail{
      \hline
    }
    \bottomcaption{#1}
  \begin{xtabular}{#2}
    #5
   \bottomrule
  \end{xtabular}
  \label{tabla:#4}
  \end{center}
}



\definecolor{cgoLight}{HTML}{EEEEEE}
\definecolor{cgoExtralight}{HTML}{FFFFFF}

%
% Nuevo comando para tablas grandes sin cabecera.
\newcommand{\tablaSinCabeceraConBandas}[5]{%
  \begin{center}
    \tablefirsthead{
      \toprule
    }
    \tablehead{
      \multicolumn{#3}{l}{\small\sl continúa desde la página anterior}\\
      \hline
    }
    \tabletail{
      \hline
      \multicolumn{#3}{r}{\small\sl continúa en la página siguiente}\\
    }
    \tablelasttail{
      \hline
    }
    \bottomcaption{#1}
    \rowcolors[]{1}{cgoExtralight}{cgoLight}

  \begin{xtabular}{#2}
    #5
   \bottomrule
  \end{xtabular}
  \label{tabla:#4}
  \end{center}
}


















\graphicspath{ {./img/} }

% Capítulos
\chapterstyle{bianchi}
\newcommand{\capitulo}[2]{
	\setcounter{chapter}{#1}
	\setcounter{section}{0}
	\chapter*{#2}
	\addcontentsline{toc}{chapter}{#2}
	\markboth{#2}{#2}
}

% Apéndices
\renewcommand{\appendixname}{Apéndice}
\renewcommand*\cftappendixname{\appendixname}

\newcommand{\apendice}[1]{
	%\renewcommand{\thechapter}{A}
	\chapter{#1}
}

\renewcommand*\cftappendixname{\appendixname\ }

% Formato de portada
\makeatletter
\usepackage{xcolor}
\newcommand{\tutor}[1]{\def\@tutor{#1}}
\newcommand{\course}[1]{\def\@course{#1}}
\definecolor{cpardoBox}{HTML}{E6E6FF}
\def\maketitle{
  \null
  \thispagestyle{empty}
  % Cabecera ----------------
\noindent\includegraphics[width=\textwidth]{cabecera}\vspace{1cm}%
  \vfill
  % Título proyecto y escudo informática ----------------
  \colorbox{cpardoBox}{%
    \begin{minipage}{.8\textwidth}
      \vspace{.5cm}\Large
      \begin{center}
      \textbf{TFG del Grado en Ingeniería Informática}\vspace{.6cm}\\
      \textbf{\LARGE\@title{}}
      \end{center}
      \vspace{.2cm}
    \end{minipage}

  }%
  \hfill\begin{minipage}{.20\textwidth}
    \includegraphics[width=\textwidth]{escudoInfor}
  \end{minipage}
  \vfill
  % Datos de alumno, curso y tutores ------------------
  \begin{center}%
  {%
    \noindent\LARGE
    Presentado por \@author{}\\ 
    en Universidad de Burgos --- \@date{}\\
    Tutor: \@tutor{}\\
  }%
  \end{center}%
  \null
  \cleardoublepage
  }
\makeatother

\newcommand{\nombre}{Daniel Puente Gabarri} %%% cambio de comando

% Datos de portada
\title{Aplicación para el Acceso a Plataformas de ELearning desde Dispositivos Móviles}
\author{\nombre}
\tutor{María Belén Vaquerizo García y Bruno Baruque Zanón}
\date{\today}

\begin{document}

\maketitle


\newpage\null\thispagestyle{empty}\newpage


%%%%%%%%%%%%%%%%%%%%%%%%%%%%%%%%%%%%%%%%%%%%%%%%%%%%%%%%%%%%%%%%%%%%%%%%%%%%%%%%%%%%%%%%
\thispagestyle{empty}


\noindent\includegraphics[width=\textwidth]{cabecera}\vspace{1cm}

\noindent D. nombre tutor, profesor del departamento de nombre departamento, área de nombre área.

\noindent Expone:

\noindent Que el alumno D. \nombre, con DNI 71347273-P, ha realizado el Trabajo final de Grado en Ingeniería Informática titulado título de TFG. 

\noindent Y que dicho trabajo ha sido realizado por el alumno bajo la dirección del que suscribe, en virtud de lo cual se autoriza su presentación y defensa.

\begin{center} %\large
En Burgos, {\large \today}
\end{center}

\vfill\vfill\vfill

% Author and supervisor
\begin{minipage}{0.45\textwidth}
\begin{flushleft} %\large
Vº. Bº. del Tutor:\\[2cm]
D. nombre tutor
\end{flushleft}
\end{minipage}
\hfill
\begin{minipage}{0.45\textwidth}
\begin{flushleft} %\large
Vº. Bº. del co-tutor:\\[2cm]
D. nombre co-tutor
\end{flushleft}
\end{minipage}
\hfill

\vfill

% para casos con solo un tutor comentar lo anterior
% y descomentar lo siguiente
%Vº. Bº. del Tutor:\\[2cm]
%D. nombre tutor


\newpage\null\thispagestyle{empty}\newpage




\frontmatter

% Abstract en castellano
\renewcommand*\abstractname{Resumen}
\begin{abstract}
En este primer apartado se hace una \textbf{breve} presentación del tema que se aborda en el proyecto.
\end{abstract}

\renewcommand*\abstractname{Descriptores}
\begin{abstract}
Palabras separadas por comas que identifiquen el contenido del proyecto Ej: servidor web, buscador de vuelos, android \ldots
\end{abstract}

\clearpage

% Abstract en inglés
\renewcommand*\abstractname{Abstract}
\begin{abstract}
A \textbf{brief} presentation of the topic addressed in the project.
\end{abstract}

\renewcommand*\abstractname{Keywords}
\begin{abstract}
keywords separated by commas.
\end{abstract}

\clearpage

% Indices
\tableofcontents

\clearpage

\listoffigures

\clearpage

\listoftables
\clearpage

\mainmatter
\capitulo{1}{Introducción}


Este trabajo de final de grado nace de la base de un proyecto anterior, debido a la necesidad de ampliar la funcionalidad de dicho proyecto para poder ser empleado también desde una aplicación móvil desarrollada en Android, que permita llevar a cabo las mismas funcionalidades y acceder a los mismos contenidos a disposición del alumno por medio del sistema E-Learning de la Universidad de Burgos.

\section{Contenido del proyecto}

\subsection{Memoria}
\begin{itemize}
	\item \textbf{Introducción:} descripción breve sobre el proyecto a desarrollar.
	\item \textbf{Objetivos del proyecto:} principales objetivos del proyecto que posteriormente derivarán el las historias de usuario.
	\item \textbf{Conceptos teóricos:} principales conceptos teóricos necesarios para la comprensión del proyecto y del producto creado.
	\item \textbf{Aspectos relevantes del desarrollo del proyecto:} principales aspectos del proyecto.
	\item \textbf{Trabajos relacionados:} principales proyectos relacionados con producto creado.
	\item \textbf{Conclusiones y Líneas de trabajo futuras:} resumen sobre los conocimientos adquiridos y aspectos de mejora posibles en un futuro.
	
\end{itemize}

\subsection{Anexos}

\begin{itemize}
	\item \textbf{Plan de Proyecto Software:} explicación sobre la planificación seguida durante el desarrollo del proyecto, junto con un estudio de viabilidad legal.
	\item \textbf{Especificación de Requisitos:} explicación de las historias de usuario, casos de uso, etc.
	\item \textbf{Especificación de diseño:} explicación y descripción del diseño de la aplicación.
	\item \textbf{Documentación técnica de programación:} explicación de los requisitos mínimos y configuración para el funcionamiento del proyecto.
	\item \textbf{Documentación de usuario:} explicación sobre el funcionamiento de la aplicación.
	
\end{itemize}

\subsection{Proyecto}

El proyecto se encuentra disponible en el siguiente enlace:

\url{https://github.com/danielpuente-dpg/GII14.K.QUICKTEST}


\capitulo{2}{Objetivos del proyecto}

A continuación, se va a llevar a cabo una citación y argumentación de los distintos objetivos a realizar en este proyecto con el objetivo de transmitir la finalidad del mismo.

\section{Objetivos funcionales}

Partiendo del proyecto de partida identificamos los siguientes nuevos objetivos.

\subsection{Docentes}

\begin{itemize}

	\item Estos cuestionarios permitirán a los docentes poder llevar a cabo una evaluación de los conocimientos adquiridos por su alumnado.
	
	\item Los profesores podrán ver la calificación de sus alumnos desde su dispositivo.
	
	\item Los docentes podrán ver el cuestionario al que se enfrentan sus alumnos.
	
	\item Todas estas acciones podrán realizarse desde el propio dispositivo del profesor.	
	
	
\end{itemize}

\subsection{Alumnos}

\begin{itemize}

	\item Podrán resolver los distintos cuestionarios a los que tenga que enfrentarse desde su dispositivo.
	
	\item Podrán ver la calificación obtenida y revisar como han resuelto el cuestionario.
	
	\item Las calificaciones de los alumnos variarán en función de las distintas recompensas.
	
	\item Los alumnos serán evaluados en función de su orden de respuesta al cuestionario.
	
	\item Los alumnos podrán utilizar comodines si lo consideran necesarios y su calificación se vera afectada en caso de uso.
	

\end{itemize}

\section{Objetivos de carácter técnico}

\begin{itemize}
	
	\item Desarrollar una aplicación Android para la resolución de estos cuestionarios que sea compatible con el mayor numero de dispositivos.
	
	\item Desarrollar un API que se capaz de proporcionar una interfaz común para la comunicación del proyecto de partida con la aplicación a implementar.
	
	\item Utilizar un sistema de control de versiones como \emph{Git} que permita registrar los cambios realizados durante el desarrollo del proyecto y la utilización de una plataforma que permita alojar este tipo de proyectos como \emph{GitHub}.
		
	\item Desarrollar una aplicación Android que cumpla el diseño Material Design.
	
	\item Utilizar el sistema de compilación Gradle en Android para: \emph{automatizar, administrar el proceso de compilación del proyecto e incluir en el mismo, ciertas configuraciones de compilación personales}.
	
	\item Realizar el proyecto mediante la utilización de la metodología ágil: \emph{Scrum}.
	
	\item Utilización del web service de Moodle para obtener la información necesaria por la aplicación.

\end{itemize}


\section{Objetivos personales}

\begin{itemize}
	
	\item Enfrentarme a nuevos retos que pongan a prueba todos mis conocimientos adquiridos a lo largo del grado.
	
	\item Aprender a desarrollar aplicaciones Android.
	
	\item Enfrentarme a un posible proyecto o reto que se asemejen a los del mundo laboral.
	
	\item Construir y comprender un API.
	
	\item Enfrentarme a un proyecto utilizando una metodología ágil.
	\item Proporcionar una nueva forma de enfrentarnos a los cuestionarios del proyecto de partida.

\end{itemize}





















\capitulo{3}{Conceptos teóricos}

En esta sección se va proceder a la explicación de ciertos conceptos teóricos necesarios para la correcta compresión de este trabajo. Además, dichos conceptos han sido necesarios para llevar a cabo la toma decisiones sobre cómo resolver el trabajo junto con la realización del mismo.


\section{Web API}

Antes de hablar de que es un Web API tendremos que explicar que es un API, ya que la funcionalidad de un Web API es similar a la de un API, pero orientada a la Web.

\subsection{API}

Una posible definición podría ser la siguiente: \emph{Un API (siglas de Application Programming Interface) es un conjunto de reglas (código) y especificaciones que las aplicaciones pueden seguir para comunicarse entre ellas: sirviendo de interfaz entre programas diferentes de la misma manera en que la interfaz de usuario facilita la interacción humano-software"} \cite{wiki:api}.





Es decir, permite la comunicación entre distintos componentes software. La mayor ventaja es que permite reutilizar métodos escrito en un determinado lenguaje o software, de esta manera evitamos la existencia de duplicidad de una misma funcionalidad en los diferentes componentes. Además, estas funcionalidades se encuentran testeadas y funcionan de forma adecuada en un determinado componente software.

\subsection{Web API}

En este caso y relacionado con lo anteriormente explicado la lógica de un Web API es la misma que un API salvo que en este caso esta comunicación, es decir, el intercambio de información se realiza entre un servicio web y una aplicación mediante una URL. Para llevar a cabo esta comunicación se utilizan peticiones HTTP o HTTPS y toda esta información se encuentra encapsulada generalmente en XML o JSON.
Existen principalmente cuatro tipos de Web API.

\subsubsection{SOAP}

Es un protocolo estándar de intercambio de información y datos en XML entre dos objetos cuyas siglas son las siguientes Simple Object Access Protocol. Posteriormente en esta misma sección se dedicará un apartado donde se explicará en mayor.

\subsection{XML-RPC}

Es un protocolo que llama a un procedimiento remoto que utiliza XML para encapsular los datos y llamadas HTTP para llevar a cabo la comunicación.

\subsubsection{JSON-RPC}

Es un protocolo cuya lógica es igual que el protocolo explicado anteriormente, salvo que en este caso utiliza el formato JSON para encapsular los datos.

\subsubsection{REST}

En este caso es una arquitectura software para sistemas hipermedia en la World Wide Web. Además, esta arquitectura utiliza el protocolo HTTP para llevar a cabo la comunicación. No obstante, en esta misma sección se dedicará un apartado para llevar a cabo una explicación en mayor detalle. \cite{wiki:rest}


\section{LTI}

\section{SOAP}

Anteriormente se ha llevado a cabo una pequeña explicación sobre este protocolo de comunicación. Relacionado con esto definíamos SOAP como un protocolo que posibilita el intercambio de información, es decir, la comunicación mediante internet entre aplicaciones. Para llevar a cabo este intercambio de datos se utiliza el formato XML, gracias a este el intercambio de datos se puede realizar independientemente de la plataforma o lenguaje utilizado.
Hoy en día, estos mensajes de tipo SOAP se envían generalmente mediante el protocolo HTTP, aunque también se pueden utilizar otros como SMTP, TCP o JMS entre otros. Otro aspecto importante es que es un pilar fundamental para los Web Services, ya que estos utilizan también el formato XML para describir los servicios. Además, permite el intercambio de datos o llamadas a procedimientos remotos mediante RPC.
Para poder llevar a cabo la correcta lectura de un mensaje SOAP el documento XML debe de tener una determinada estructura. Dicha estructura debe estar formada por las siguientes partes:

\begin{itemize}

	\item \textbf{Envelope:} esta parte de la estructura del mensaje es obligatoria, ya que identifica el mensaje.
	
	\item \textbf{Header:} esta parte de la estructura permite enviar información adicional sobre cómo debe de procesarse el mensaje. Es por esto por lo que esta parte no es obligatoria, debido a que únicamente aporta cierta información relacionada con la lectura del mismo.
	
	\item \textbf{Body:} esta parte de la estructura es la encargada de almacenar toda la información del mensaje.
	
	\item \textbf{Fault:} esta parte de la estructura al igual que el Header es opcional, ya que aporte información relacionada con ciertos errores producidos durante el procesado del mensaje.
	
\end{itemize}


\imagen{EstructuraSoap}{Estructura de un mensaje SOAP \cite{wiki:soap} \cite{wiki:soap2}}

\section{REST}

Como ya hemos explicado anteriormente, es un estilo de arquitectura software para sistemas hipermedia distribuidos como la World Wide Web. \cite{wiki:rest2} cuyas siglas son las siguientes REpresentational State Transfer o en castellano Transferencia de Estado Representacional.
En la actualidad, el termino REST se utiliza para aquellas interfaces que utilicen HTTP para el intercambio de información entre sistemas pudiendo utilizar cualquier formato para encapsular los datos, aunque generalmente los más utilizados son XML y JSON. Además, actualmente los sistemas que siguen las pautas o principios REST también se les suelen denominar RESTful.
Para que un sistema se considere RESTful debe de cumplir con las siguientes pautas o principios.

\begin{itemize}

	\item \textbf{Protocolo cliente/servidor sin estado:} de manera que cada petición HTTP debe contener toda aquella información necesaria para poder ejecutarse, es decir, de esta forma evitamos que tanto el cliente como el servidor tengan que almacenar información sobre el estado previo de la misma para poder llevarla a cabo.
	
	\item \textbf{Operaciones:} un sistema REST debe poder llevar cabo las siguientes operaciones: POST (crear), GET (leer o consultar), PUT (editar) y DELETE (borrar). Estas operaciones como podemos observar se asemejan en gran parte a las operaciones CRUD en bases de datos.
	
	\item \textbf{Sistema de capas:} el sistema deberá de utilizar una arquitectura jerárquica entre los distintos componentes que la formen, de esta manera garantizamos que cada una de estas capas se encargue de llevar a cabo una única funcionalidad.
	
	\item \textbf{Manipulación de recursos:} para realizar la manipulación de los objetos se lleva a cabo mediante la URI. Dicha URI se utiliza como el identificador único para cada recurso, un recurso es un elemento de información. De esta manera se simplifica el acceso a la información para su posterior manipulación.
	
	\item \textbf{Uso de hipermedios:} para las transiciones entre los distintos estados de la aplicación y para la información de la misma.
	
\end{itemize}

\cite{wiki:rest}
 

\capitulo{4}{Técnicas y herramientas}

En esta sección se llevará a cabo una mención y breve explicación sobre el conjunto técnicas y herramientas utilizadas durante el desarrollo del proyecto.

\section{Lenguajes}

\subsection{PHP}

PHP es un lenguaje de código abierto muy popular especialmente adecuado para el desarrollo web y que puede ser incrustado en HTML \cite{wiki:php} cuyas siglas es un acrónimo recursivo de PHP Hypertext Preprocesssor. Este lenguaje se utiliza principalmente en el lado del servidor y está enfocado en el desarrollo web para el contenido dinámico \cite{wiki:php2}.
Este lenguaje se ha decidido utilizar frente a otras alternativas en el lado del servidor debido principalmente a que el proyecto de partida se basa en este lenguaje para el acceso a la base de datos. Además, cabe destacar que desde mi punto de vista me parece una gran decisión al ser uno de los lenguajes más utilizados en el lado del servidor, existe una gran documentación al respecto y la curva de aprendizaje es menos costosa al asemejarse a otros lenguajes orientados a objetos.

\subsection{JSON}\label{json}

JSON es un formato ligero de intercambio de datos cuyas siglas es un acrónimo de JavaScript Object Notation, es decir, Notación de Objetos de JavaScript \cite{wiki:json}. JSON es un formato de texto independientemente del lenguaje y se le considera una gran alternativa frente al formato de intercambio de datos XML al ser mucho más sencillo de parsear por un analizador sintáctico \cite{wiki:json2}.
Este lenguaje se ha decidido utilizar para la devolución de los diferentes datos a la aplicación Android. De esta forma, la comunicación de la aplicación se podrá realizar de una forma más sencilla con la base de datos. Además, las propias librerías de Android contienen ciertas clases para poder realizar de forma fácil un correcto tratamiento de los datos.

\subsection{Android}

\section{Herramientas}

\subsection{XAMPP}

XAMPP es un paquete de instalación independiente de la plataforma y de software libre. Dicho paquete contiene el sistema de gestión de bases de datos MySQL¸ el servidor web Apache y los interpretes para lenguajes de script: PHP y Perl \cite{wiki:xampp}. Su nombre proviene del acrónimo X indicando que es compatible para cualquier sistema operativo y el resto de siglas AMPP hace referencia a cada uno de los elementos que contiene el paquete, es decir, Apache, MySQL, PHP, Perl respectivamente.
Esta herramienta se ha decidido utilizar principalmente al ser utilizada en el proyecto de partida. Además, cabe destacar que me parece una gran decisión ya que el propio paquete contiene todo lo necesario para poder crear un servidor web.

\subsection{GitHub}

GitHub es una plataforma de desarrollo colaborativo para alojar proyectos utilizando un sistema de control de versiones Git \cite{wiki:github}. 
Se ha decidido optar por esta herramienta al haberla utilizado en ciertas ocasiones, al ser una de las más utilizadas para el desarrollo colaborativo y al existir un plugin ZenHub para la gestión desde el propio repositorio de las distintas tareas a realizar en un panel de trabajo. Además, cabe destacar que el repositorio permite dos tipos de formato para el repositorio: público y privado. En nuestro caso se encuentra público. \url{https://github.com/danielpuente-dpg/GII14.K.QUICKTEST}

\subsection{ZenHub}\label{zenhub}

\subsection{TortoiseSVN}

TortoiseSVN es un software libre utilizado para llevar a cabo el control de versiones del proyecto sobre GitHub \cite{wiki:tortoisesvn}. Implementa una extensión del shell de Windows, por lo que el control de versiones no tiene por qué realizarse mediante línea de comandos, sino que se puede realizar mediante la interfaz gráfica que proporciona esta herramienta la cual es muy fácil de utilizar.

\subsection{Moodle}
\subsection{PhpStorm}
\subsection{phpMyAdmin}
\subsection{PHPUnit}
\subsection{SonarQube}
\subsection{Advanced REST Client}
\subsection{StarUML}
\subsection{Librerías}

A continuación, se incluye una pequeña explicación sobre las librerías utilizadas para el desarrollo del proyecto.

\subsubsection{Consumir peticiones}

Durante el desarrollo de la aplicación se han barajado ambas librerías para en consumo de las peticiones al propio API y al webservice de Moodle. Finalmente, se ha decido optar por la librería Retrofit al ser más fácil de utilizar, más intuitiva a la hora de crear y configurar las distintas peticiones, más flexible y suele tener mejores resultados que Volley

\begin{itemize}

	\item \textbf{Retrofit:} Retrofit proporciona un framework para poder interactuar con APIs y enviar peticiones HTTP de forma fiable desde aplicaciones Android. Para ello proporciona un cliente HTTP encargado de interactuar con APIs y de gestionar las peticiones de manera transparente.
Esta librería es utilizada para el tratamiento de las peticiones desde la aplicación android. Además, permite consumir peticiones de manera síncrona y asíncrona. \cite{wiki:retrofit}
	
	\item \textbf{Volley:} Al igual que la librería anterior, Volley es una librería desarrollada por Google para enviar peticiones HTTP desde aplicaciones Android. \cite{wiki:volley}
	
	\item \textbf{Gson:} Esta librería proporcionada por Google, permite trabajar con JSON a la hora de serializar y deserializar los objetos. Es utilizada para convertir las respuestas JSON de las APIs en objetos Java de manera muy sencilla, eliminado toda esta carga de conversión al programador. Además, también permite el paso contrario a la hora de enviar objetos a las APIs, al encargarse de transformar estos objetos en respuestas en formato JSON. \cite{wiki:gson}

\end{itemize}

\section{Técnicas}

\subsection{LTI}
\subsection{MVC}
\subsection{SCRUM}\label{scrum} 

SCRUM es un modelo de referencia que define un conjunto de práctica y roles, que pueden utilizarse como punto de partida para llevar a cabo la definición de como se elaborará el proceso de desarrollo del proyecto \cite{wiki:scrum}. Cabe destacar que este modelo se encuentra dentro de los marcos de las metodologías ágiles y que es de los más utilizados actualmente.
Los roles principales son:

\begin{itemize}

	\item \textbf{Scrum Master:} se encarga de gestionar los cambios y procura facilitar la aplicación de esta metodología.
	
	\item \textbf{Product Owner:} en este grupo se encontrarán el personal interno o externo que representa al cliente. Este grupo se encargará de que el trabajo se realice de forma acorde con las necesidades y peticiones del cliente.
	
	\item \textbf{Team:} representa al equipo de desarrollo encargados de ejecutar el desarrollo y de entregar el producto deseado.
	
	\item \textbf{Product Owner:} en este grupo se encontrarán el personal interno o externo que representa al cliente. Este grupo se encargará de que el trabajo se realice de forma acorde con las necesidades y peticiones del cliente.	

\end{itemize}
\cite{wiki:scrum}
 
Para realizar el desarrollo del producto se deben de definir un Product Backlog, el cual contendrá todas las historias de usuario a realizar. Dichas historias de usuario se asignarán a los distintos Sprints o iteraciones que se realicen a lo largo del desarrollo del producto hasta finalizar con todas las historias de usuario. 











\capitulo{5}{Aspectos relevantes del desarrollo del proyecto}

Este apartado pretende recoger los aspectos más interesantes del desarrollo del proyecto, comentados por los autores del mismo.
Debe incluir desde la exposición del ciclo de vida utilizado, hasta los detalles de mayor relevancia de las fases de análisis, diseño e implementación.
Se busca que no sea una mera operación de copiar y pegar diagramas y extractos del código fuente, sino que realmente se justifiquen los caminos de solución que se han tomado, especialmente aquellos que no sean triviales.
Puede ser el lugar más adecuado para documentar los aspectos más interesantes del diseño y de la implementación, con un mayor hincapié en aspectos tales como el tipo de arquitectura elegido, los índices de las tablas de la base de datos, normalización y desnormalización, distribución en ficheros3, reglas de negocio dentro de las bases de datos (EDVHV GH GDWRV DFWLYDV), aspectos de desarrollo relacionados con el WWW...
Este apartado, debe convertirse en el resumen de la experiencia práctica del proyecto, y por sí mismo justifica que la memoria se convierta en un documento útil, fuente de referencia para los autores, los tutores y futuros alumnos.

\capitulo{6}{Trabajos relacionados}

Durante el desarrollo del proyecto se analizó el comportamiento y diseño de ciertas aplicaciones. No obstante, cabe destacar las siguientes.

\section{Proyecto de partida}

Como se ha comentado anteriormente, este proyecto nace por la necesidad de ampliar la funcionalidad de la resolución de este tipo de cuestionarios desde nuestro propio dispositivo Android, ya que presenta cierta problemática en el renderizado de la aplicación en un dispositivo móvil.

Al tratarse de una ampliación de este proyecto tanto las ventajas como inconvenientes que presenta el proyecto de partida serán los mismo que la aplicación desarrollada.


\section{Cuestionarios de Moodle}

El propio entorno educativo de Moodle incorpora sus propios cuestionarios. Destacar que la herramienta a diferencia de la desarrollada es mucho más completa. 

\subsection{Desventajas}

A continuación, se citan algunas de las principales ventajas que presentan estos cuestionarios nativos de Moodle frente al producto creado.

\begin{itemize}
	\item \textbf{Multirespuesta:} permite que en el número de respuestas de una pregunta sea mayor que uno.
	
	\item \textbf{Categorías:} permite el almacenamiento de las preguntas por categorías.
	
	\item \textbf{Reintentos:} el profesor puede configurar el número de intentos de un mismo cuestionario.
	
	\item \textbf{Tiempo:} permite incluir un tiempo en el que el cuestionario tiene que resolverse.
	
	\item \textbf{Imagenes:} permite incluir en la propia pregunta una image.
	
	\item \textbf{Aleatoriedad:} permite que en un mismo cuestionario el número de preguntas varíen entre alumnos o intentos.
	
\end{itemize}

\subsection{Ventajas}

A diferencia de los propios cuestionarios de Moodle la aplicación permite:

\begin{itemize}
	\item \textbf{Dinamismo:} la herramienta construida permite que los alumnos se enfrenten a los cuestionarios de una manera más divertida, debido al sistema de recompensa empleado: \emph{comodines y orden.}
	
	\item \textbf{Comodines y orden:} los alumnos se enfrentarán a estos cuestionarios no solamente con la motivación de superarlo, sino que, \emph{si son los primeros obtendrán una recompensa}, \emph{si son de los últimos en responder tendrán a su disposición un mayor número de comodines}, es decir \emph{fomentará la competición entre los alumnos}, \emph{ayudará a que los propios alumnos se relacionen fruto de la competición}, etc.
	
	\item \textbf{Móvil:} al permitir que los usuarios puedan resolver estos cuestionarios desde su dispositivo Android, estos cuestionarios podrán realizarse como una comprobación al finalizar cada clase de lo aprendido en la misma, al necesitar únicamente de sus teléfonos móviles.
\end{itemize}

\section{Big Web Quiz}\cite{wiki:bigwebquiz}

Durante el desarrollo del proyecto se realizo una búsqueda de alguna herramienta de resolución de cuestionarios en dispositivos Android como complemento de Moodle pero, no se encontró ninguna al respecto. No obstante, se ha elegido esta herramienta: \emph{Big Web Quiz}al ser desarrollada por Google y al fomentar, al igual que la aplicación construida: \emph{\textbf{la competición entre los participantes.}}

\imagen{bigwebquiz}{Big Web Quiz en el Play Store.}

Ambas aplicaciones destacan por:
\begin{itemize}
	\item Fomentar la competición.
	\item Proporcionan una forma más dinámica de enfrentarse a un cuestionario.
	\item Permiten resolver los cuestionarios desde un dispositivo móvil Android.
	\item Fomentar la gamificacion durante la resolución de los cuestionaros, es decir aprender divirtiéndose.
\end{itemize}

\subsection{Desventajas}

A diferencia de la aplicación desarrollada, Big Web Quiz proporciona:
	\begin{itemize}
	\item Permite varias formas de respuesta: multirespuesta, preguntas con imágenes, proporciona gráficos, entre otras.
	
	\item Permite la selección de avatares al iniciar un cuestionario.
	
	\item Fomenta la competición de manera agrupada en varios cuestionarios.
	
	\end{itemize}
	
	

\subsection{Ventajas}

\begin{itemize}
	\item \textbf{Moodle:} la aplicación desarrollada es compatible con LTI y por tanto con Moodle. Sin embargo, Big Web Quiz no. \ref{subsec:LTI}.
	
	\item \textbf{Docentes:} la aplicación desarrollada permite a los docentes obtener la calificación de sus alumnos.
	
	\item Para su funcionamiento no es necesario un Chromecast.
	
	\item Big Web Quiz solamente permite enfrentarse a un cuestionario desde un dispositivo móvil.
	
	\item La aplicación desarrollada permite que los alumnos puedan obtener información sobre como han resuelto un cuestionario.
\end{itemize}





\capitulo{7}{Conclusiones y Líneas de trabajo futuras}

En esta sección se incluyen todas las conclusiones derivadas del desarrollo del proyecto y las posibles lineas de trabajo futuras.

\section{Conclusiones}

\begin{itemize}
	\item Desde mi punto de vista, considero que los objetivos principales del proyecto han sido superados durante el desarrollo del mismo.
	
	\item Al haber tenido que basarse en un proyecto ya desarrollado, he tenido muchas ventajas al tener muchas funcionalidades necesarias ya implementadas, pero también he tenido el inconveniente de que al crear nuevas funcionalidades estas no alteren el funcionamiento inicial.
	
	\item La creación de un API sobre un proyecto web ha aportado no solamente nuevos conocimientos sobre: \emph{PHP, lado del servidor, JavaScript, etc}. Sino, que me ha permitido comprender que prácticamente todas las aplicaciones Android de productos que anteriormente habían sido desarrollados para la web, han tenido que crear un API para manejar esta lógica.
	
	\item Se ha obtenido un nuevo conocimiento sobre la programación en Android, PHP o los web services que hasta la realización de este se desconocían.
	
	\item Se ha comprendido el verdadero potencial de la utilización de herramientas para: \emph{el control de versiones, la utilización de una metodología ágil, utilización de sistemas de construcción basado en dependencias como Gradle, herramienta de integración continua como SonarQube, etc}.
	
	\item Se ha comprendido como puede ser la dinámica en el mundo laboral en un futuro, es decir la asignación de un proyecto, su gestión, nuevos conocimientos a adquirir, metodología de trabajo, etc.
\end{itemize}



\section{Lineas de trabajo futuras}

\begin{itemize}
	\item \textbf{Proyecto de partida:}
	\begin{list}{-}
		\item Modificar la clase encargada de manejar la conexión a la base de datos implementando un patrón Singleton. De esta forma, limitaremos a que exista un único hilo para acceder a la misma.
		\item Modificar todas aquellas consultas a la base de datos que no se realicen mediante sentencias preparadas. De esta manera, evitaremos posibles problemas de seguridad y mejoraremos el mantenimiento.
		\item Mejorar la forma de crear un cuestionario, ya que el método encargado de esta lógica se ocupa de tareas: \emph{crear y editar un cuestionario.}
		
		\item Incorporar un sistema que permita implementar la concurrencia.
	\end{list}
	
	\item \textbf{Aplicación Android:}
	\begin{list}{-}
		\item Permitir que la calificación de los alumnos pueda ser almacenada en el apartado de calificaciones del curso,  mediante alguna biblioteca en Android que permita el envío de datos a Moodle cumpliendo el estándar LTI.
		
		\item Posibilite al usuario una retroalimentación sobre las respuestas realizadas y la respuesta correcta.
		
		\item Permita la creación de cuestionarios multirespuesta o que las preguntas puedan adjuntar una imagen.
		
		\item Nuevas formas de enfrentarse a un cuestionario.
		
		\item Los docentes puedan obtener ciertos gráficos sobre: \emph{las respuestas más contestadas, aquellas preguntas con derecho a comodín, entre otras.}
		\item Permita que los docentes puedan crear, modificar, duplicar o eliminar sus cuestionarios. Destacar que los métodos principales para desarrollar esta lógica, se encuentran implementados en el propio API.
		\item Realizar una internacionalización de la aplicación.
		\item La aplicación incluya soporte para personas con dificultades como: \emph{problemas de vista o auditivas.} Algunos ejemplos podrían ser: \emph{asistente de lectura de voz o paleta de colores propia para personas con problemas visuales.}
	\end{list}
\end{itemize}



\bibliographystyle{plain}
\bibliography{bibliografia}

\end{document}
